\section{Methodology}\label{sec:methodology}

This section describes the main relations that make up the \textsc{Terrae}
proof system: how the forest is committed, how inference is checked, how forest
well-formedness is certified, and how the training proof reduces split
selection to algebraic relations over histograms.

\subsection{Forest Representation}\label{subsec:method-forest-representation}

We adopt the hyperbox forest encoding of \cite{cq++24}. Each node carries a
feature-wise half-open interval box, and a sample reaches a leaf exactly when
its quantized feature vector lies inside that box. 
We assume there are $N_{\mathrm{tot}}$ nodes, and
all committed objects share
one global node namespace, so node identifiers are
$u \in \{0,\ldots,N_{\mathrm{tot}}-1\}$ everywhere. Node $0$ is a dummy node, and
all of its committed attributes are fixed to the dummy value $0$. When a node identifier labels
a committed row, the corresponding physical row in the row-major encoding is
$u$. The committed
forest consists of the commitments of:
\begin{itemize}
  \item left and right child index vectors
  $\mathbf{lc},\mathbf{rc}\in[N_{\mathrm{tot}}]^{N_{\mathrm{tot}}}$,
  and parent index vector $\mathbf{par}\in[N_{\mathrm{tot}}]^{N_{\mathrm{tot}}}$, where invalid entries are set $0$;
  \item lower and upper node bounds
  $\mathbf{N}^{\leftarrow},\mathbf{N}^{\rightarrow}\in[B+1]^{N_{\mathrm{tot}}\times d}$,
  where row $u$ stores the $d$-dimensional half-open box of node $u$ and every
  leaf row describes one accepting region of the tree;
  \item a feature-selection matrix
  $\mathbf{E}\in\{0,1\}^{N_{\mathrm{tot}}\times d}$,
  whose internal-node rows are one-hot and indicate which feature is split at
  each node, while leaf rows are zero;
  \item a threshold-index vector
  $\boldsymbol{\Theta}\in\{0,1,\ldots,B\}^{N_{\mathrm{tot}}}$,
  whose internal-node entries record the selected threshold bin and whose leaf rows are zero;
  \item a tree-index vector $\mathbf{tree}\in[TK+1]^{N_{\mathrm{tot}}}$, where tree identifier $0$ is the dummy value and real trees are ranked from $1$ to $TK$, and a
  root indicator $\mathbf{root}\in\{0,1\}^{N_{\mathrm{tot}}}$;
  \item a depth indicator $\mathbf{depth}\in\{0,1,\ldots,h\}^{N_{\mathrm{tot}}}$, where $h$ is the maximum depth of the trees in the forest. This is used for proving the correctness of tree construction.
  \item a leaf-score vector $\mathbf{score}\in\mathbb{F}^{N_{\mathrm{tot}}}$,
  whose leaf entries store the learning-rate-scaled scalar contribution
  returned by that tree and whose internal-node entries are zero;
  \item a boolean vector $\mathbf{is\_leaf}\in\{0,1\}^{N_{\mathrm{tot}}}$, 
  where $\mathbf{is\_leaf}[u]=1$ if and only if $u$ is a leaf node;
  \item a boolean vector $\mathbf{valid}\in \{0, 1\}^{N_{\mathrm{tot}}}$, where $\mathbf{valid}[u]=0$ iff the node is dummy.
\end{itemize}

We assume there are $N_{\mathrm{int}}$ internal nodes and $N_{\mathrm{leaf}}$ leaf nodes
in the forest, where internal nodes are indexed by $\{1,\ldots,N_{\mathrm{int}}\}$ and leaf nodes are indexed by
$\{N_{\mathrm{int}}+1, \dots, N_{\mathrm{tot}}-1\}$.

Note that $\mathbf{N}^{\leftarrow}$, $\mathbf{N}^{\rightarrow}$, $\mathbf{lc}$
and $\mathbf{rc}$ together fix the entire forest.
Other vectors/matrices are auxiliary to the proof
of training and inference.

\Cref{fig:forest-encoding-inference} shows the intuition behind
the forest encoding: the left panel illustrates the structure of the forest,
while the right panel presents the committed objects.

\begin{figure*}[t]
  \centering
  \resizebox{\linewidth}{!}{%
  \begin{tikzpicture}[
    >=Latex,
    font=\small,
    paneltitle/.style={font=\bfseries},
    treelabel/.style={font=\bfseries\small},
    internal/.style={
      circle,
      draw=blue!55!black,
      line width=0.55pt,
      fill=blue!5,
      minimum size=10.2mm,
      inner sep=1pt,
      align=center
    },
    leaf/.style={
      rounded corners=2pt,
      draw=green!50!black,
      line width=0.55pt,
      fill=green!10,
      minimum width=14mm,
      minimum height=8.6mm,
      inner sep=2pt,
      align=center
    },
    note/.style={align=center, font=\scriptsize},
    bigtable/.style={
      matrix of nodes,
      ampersand replacement=\&,
      row sep=0.75mm,
      column sep=\colgap,
      nodes={
        anchor=center,
        inner xsep=2.5pt,
        inner ysep=1.2pt,
        minimum height=4.5mm,
        text depth=.25ex,
        outer sep=0pt
      },
      row 1/.style={nodes={font=\bfseries}},
      column 1/.style={nodes={minimum width=7mm}},
      column 2/.style={nodes={minimum width=11mm}},
      column 3/.style={nodes={minimum width=10mm}},
      column 4/.style={nodes={minimum width=13mm}},
      column 5/.style={nodes={minimum width=9mm}},
      column 6/.style={nodes={minimum width=13mm}},
      column 7/.style={nodes={minimum width=13mm}},
      column 8/.style={nodes={minimum width=12mm}}
    },
    colbg/.style={
      draw=gray!70,
      rounded corners=2pt,
      line width=.5pt,
      fill=white,
      inner sep=2.8pt
    },
    meta/.style={
      draw=gray!70,
      rounded corners=2pt,
      fill=gray!6,
      align=center,
      inner sep=6pt
    },
    fetchrow/.style={
      draw=red!75!black,
      rounded corners=2pt,
      line width=.85pt,
      inner xsep=4pt,
      inner ysep=2pt
    },
    fetchtxt/.style={text=red!75!black, font=\bfseries},
    outer/.style={
      draw=gray!55,
      dashed,
      rounded corners=4pt,
      inner sep=8pt
    }
  ]

    \def\colgap{4.2mm}

    % =========================
    % Left panel: tree semantics
    % =========================
    \node[paneltitle] (ltitle) at (-4.8,3.12) {(a) Tree semantics};

    % Tree 1
    \node[treelabel] at (-7.55,2.42) {tree $1$};
    \node[internal] (t0r) at (-7.55,1.72) {$1$\\[-1mm]\scriptsize$x_1<3$};
    \node[internal] (t0l) at (-9.05,0.38) {$2$\\[-1mm]\scriptsize$x_2<4$};
    \node[leaf]     (t0rleaf) at (-6.00,0.38) {$7$\\[-1mm]\scriptsize$w_2$};
    \node[leaf]     (t0w0) at (-10.00,-1.02) {$5$\\[-1mm]\scriptsize$w_0$};
    \node[leaf]     (t0w1) at (-8.35,-1.02) {$6$\\[-1mm]\scriptsize$w_1$};

    \draw[->, line width=.75pt] (t0r) -- (t0l);
    \draw[->, line width=.75pt] (t0r) -- (t0rleaf);
    \draw[->, line width=.75pt] (t0l) -- (t0w0);
    \draw[->, line width=.75pt] (t0l) -- (t0w1);

    % Tree 2
    \node[treelabel] at (-2.65,2.42) {tree $2$};
    \node[internal] (t1r) at (-2.65,1.72) {$3$\\[-1mm]\scriptsize$x_1<2$};
    \node[leaf]     (t1lleaf) at (-4.15,0.38) {$8$\\[-1mm]\scriptsize$w'_0$};
    \node[internal] (t1rr) at (-1.15,0.38) {$4$\\[-1mm]\scriptsize$x_2<6$};
    \node[leaf]     (t1w1) at (-2.05,-1.02) {$9$\\[-1mm]\scriptsize$w'_1$};
    \node[leaf]     (t1w2) at (-0.40,-1.02) {$10$\\[-1mm]\scriptsize$w'_2$};

    \draw[->, line width=.75pt] (t1r) -- (t1lleaf);
    \draw[->, line width=.75pt] (t1r) -- (t1rr);
    \draw[->, line width=.75pt] (t1rr) -- (t1w1);
    \draw[->, line width=.75pt] (t1rr) -- (t1w2);

    % =====================================
    % Right panel: one unified committed table
    % =====================================
    \node[paneltitle] (rtitle) at (7.25,3.12) {(b) Committed row view used by the proof};

    \matrix[bigtable, anchor=north west] (tbl) at (1.15,2.55) {
      $\mathbf{u}$ \& $\mathbf{tree}$ \& $\mathbf{is\_leaf}$ \& $\mathbf{E}$ \& $\boldsymbol{\Theta}$ \& $\mathbf{N}^{\leftarrow}$ \& $\mathbf{N}^{\rightarrow}$ \& $\mathbf{score}$ \\
      $0$ \& $0$ \& $0$ \& $(0,0)$ \& $0$ \& $(0,0)$ \& $(0,0)$ \& $0$ \\
      $1$ \& $1$ \& $0$ \& $(1,0)$ \& $3$ \& $(0,0)$ \& $(B,B)$ \& $0$ \\
      $2$ \& $1$ \& $0$ \& $(0,1)$ \& $4$ \& $(0,0)$ \& $(3,B)$ \& $0$ \\
      $3$ \& $2$ \& $0$ \& $(1,0)$ \& $2$ \& $(0,0)$ \& $(B,B)$ \& $0$ \\
      $4$ \& $2$ \& $0$ \& $(0,1)$ \& $6$ \& $(2,0)$ \& $(B,B)$ \& $0$ \\
      $5$ \& $1$ \& $1$ \& $(0,0)$ \& $0$ \& $(0,0)$ \& $(3,4)$ \& $w_0$ \\
      $6$ \& $1$ \& $1$ \& $(0,0)$ \& $0$ \& $(0,4)$ \& $(3,B)$ \& $w_1$ \\
      |[fetchtxt]| $7$ \& |[fetchtxt]| $1$ \& |[fetchtxt]| $1$ \& |[fetchtxt]| $(0,0)$ \& |[fetchtxt]| $0$ \& |[fetchtxt]| $(3,0)$ \& |[fetchtxt]| $(B,B)$ \& |[fetchtxt]| $w_2$ \\
      $8$ \& $2$ \& $1$ \& $(0,0)$ \& $0$ \& $(0,0)$ \& $(2,B)$ \& $w'_0$ \\
      |[fetchtxt]| $9$ \& |[fetchtxt]| $2$ \& |[fetchtxt]| $1$ \& |[fetchtxt]| $(0,0)$ \& |[fetchtxt]| $0$
        \& |[fetchtxt]| $(2,0)$ \& |[fetchtxt]| $(B,6)$ \& |[fetchtxt]| $w'_1$ \\
      $10$ \& $2$ \& $1$ \& $(0,0)$ \& $0$ \& $(2,6)$ \& $(B,B)$ \& $w'_2$ \\
    };

    % Column background boxes
    \begin{scope}[on background layer]
      \foreach \c in {1,...,8}{
        \node[colbg, fit=(tbl-1-\c)(tbl-12-\c)] {};
      }
    \end{scope}

    % Horizontal separators: after header and after internal rows
    \foreach \r in {1,6}{
      \draw[gray!60, line width=.4pt]
        ([xshift=-2pt]tbl-\r-1.south west) -- ([xshift=2pt]tbl-\r-8.south east);
    }

    % Highlight fetched rows (one row per tree)
    \node[fetchrow, fit=(tbl-9-1)(tbl-9-8)] (fetchA) {};
    \node[fetchrow, fit=(tbl-11-1)(tbl-11-8)] (fetchB) {};

    % Query box
    \node[meta, anchor=north west, text width=3.35cm] (sample)
      at ([xshift=6mm]tbl.north east)
      {query $\mathbf{x}=(4,2)$\\[1mm]
       fetch one row per tree:\\
        $u_1=7,\;u_2=9$\\[1mm]
       and prove, for any fetched row $u$,\\
       $\mathbf{N}^{\leftarrow}[u]\le \mathbf{x}<\mathbf{N}^{\rightarrow}[u]$};

    % Arrows to the two fetched rows
    \draw[->, very thick, red!75!black]
      ([yshift=7pt]sample.west) to[out=180,in=0] ([xshift=2pt]fetchA.east);
    \draw[->, very thick, red!75!black]
      ([yshift=-7pt]sample.west) to[out=180,in=0] ([xshift=2pt]fetchB.east);

    % Outer frame
    \node[
      outer,
      fit=(ltitle)(rtitle)
          (t0r)(t0l)(t0rleaf)(t0w0)(t0w1)
          (t1r)(t1lleaf)(t1rr)(t1w1)(t1w2)
          (tbl)(sample)
    ] {};

  \end{tikzpicture}%
  }
  \caption{An illustration of the forest encoding and inference semantics.
  The left panel shows two compact trees with
  global node identifiers, and the right panel shows the committed
  objects used by inference. The highlighted entries are the fetched rows for one
  query, one per tree. Batched inference performs this fetch for every
  tree-sample pair in one shot and then aggregates scores across trees for
  each sample-class pair.}
  \Description{Two example decision trees and their committed table
  representation. The table includes node identifiers, tree identifiers,
  leaf indicators, split-feature masks, thresholds, hyper-box bounds, and leaf
  scores. Two leaf rows are highlighted to show how inference fetches one row
  per tree and checks interval containment for a query point.}
  \label{fig:forest-encoding-inference}
\end{figure*}


\subsection{Proof of Inference}\label{subsec:method-inference}

To prove the correctness of inference over batched input $\mathbf{I}\in\mathbb{F}^{b\times d}$,
the prover commits $\mathbf{U}^{\leftarrow}, \mathbf{U}^{\rightarrow}\in \mathbb{F}^{TKb\times d}$,
$\mathbf{sc}\in\mathbb{F}^{TKb}$,
and uses table lookup to prove
\begin{align}
  \mathbf{1} \| \mathbf{ind\_tree}  \| \mathbf{U} ^{\leftarrow } \| \mathbf{U} ^{\rightarrow } \| \mathbf{sc} \subseteq
  \mathbf{is\_leaf} \| \mathbf{tree} \| \mathbf{N} ^{\leftarrow }\| \mathbf{N} ^{\rightarrow }\| \mathbf{score},\label{eq:inference-lookup-1}
\end{align}
where notation ``$\subseteq$'' means that each row of the left matrix equals some row of the right matrix, and
\[
\mathbf{ind\_tree} \coloneqq (\underbrace{1, \dots, 1}_b, \underbrace{2, \dots, 2}_{b}, \dots, \underbrace{TK, \dots, TK}_b)^\top.
\]
Equivalently, the logical tree pair $(t,k)$ denotes the committed tree whose
identifier is $tK+k+1$; identifier $0$ is reserved for dummy rows.
As a sanity check,
the LHS of \Cref{eq:inference-lookup-1} has shape $TKb\times (2d+3)$,
and the RHS has shape $N_{\mathrm{tot}}\times (2d+3)$, so the table lookup is well-defined.
To prove this lookup, we pad each column vector to a $d$ column matrix,
randomly combine the matrices with a verifier challenge,
and use the table lookup of~\cite{cq++24} to prove the lookup
with (online) prover time $O(TKb)$.

This lookup argument guarantees that each row of 
$\mathbf{U}^{\leftarrow}, \mathbf{U}^{\rightarrow}$
and $\mathbf{sc}$ are aligned for the same leaf.
By further proving
\begin{align}
  \mathbf{U} ^{\leftarrow }_{\{T\times K\times b\times d\}}[t, k, i, j] \leq \mathbf{I}[i,j] < \mathbf{U}_{\{T\times K\times b\times d\}} ^{\rightarrow }[t, k, i, j], \label{eq:inference-interval}
\end{align}
the prover proves that the leaf reached by each sample is correct.
What remains is to prove the result of inference from $\mathbf{sc}$,
where $\mathbf{sc}_{\{T\times K\times b\}}[t,k,i]$ represents the score sample $i$
received from the tree trained at round $t$ for class $k$.
This includes proving the sum of $\mathbf{sc}_{\{T\times K\times b\}}$
along the second dimension and proving the location of the maximum.

\paragraph{Complexity.}
The dominant inference relation has one lookup row for every
tree-sample pair and an interval check for every feature of that row. Hence the
online prover time is
\[
    \tilde{O}(TKbd),
\]
with lower-order terms for score aggregation and the final maximum over the
$K$ classes. After the lookup, interval, sum, and maximum checks are batched,
the verifier opens only a constant number of committed polynomials at
logarithmically many points, so the verifier time and proof size are
$\tilde{O}(\log(TKbd))$ field/group elements under the KZG backend.

\subsection{Proof of Forest Well-formedness}\label{subsec:method-well-formedness}\label{subsec:proof of well-formedness}

First, the prover proves that the committed forest is well-formed, by showing
\begin{enumerate}
  \item Consistency between indicators, i.e. for any node $i$,
  \begin{align}
    \mathbf{valid}[i]=1 &\iff \mathbf{depth}[i] \neq 0,\\
    \mathbf{valid}[i]=1 &\iff \mathbf{tree}[i] \neq 0,\\
    \mathbf{valid}[i]=0 &\implies \mathbf{is\_leaf}[i] = 0.
  \end{align}

  \item Indicators are boolean, i.e.
  \begin{align}
    \mathbf{root}\odot(\mathbf{root}-\mathbf{1}) = \mathbf{0},\\
    \mathbf{is\_leaf}\odot(\mathbf{is\_leaf}-\mathbf{1}) = \mathbf{0},\\
    \mathbf{valid}\odot(\mathbf{valid}-\mathbf{1}) = \mathbf{0}.
  \end{align}

  \item No ghost nodes, i.e. all invalid nodes $u$ with $\mathbf{valid}[u]=0$ are dummy, with
    \begin{align}
    \mathbf{lc}[u] = \mathbf{rc}[u] &= 0,\\
    \forall j\in [d].\ \mathbf{N}^{\leftarrow}[u, j] = \mathbf{N}^{\rightarrow}[u, j] &= 0,\\
    \mathbf{tree}[u] = \mathbf{root}[u] = \mathbf{is\_leaf}[u] &= 0.
    \end{align}

  \item There are in total $TK$ trees, i.e.
  \begin{align}
    \sum_{u=0}^{N_{\mathrm{tot}}-1} \mathbf{root}[u] = TK.
  \end{align}

  \item The depth of root is $1$, i.e.
  \begin{align}
    \mathbf{root}\odot (\mathbf{depth}-\mathbf{1}) = \mathbf{0}.
  \end{align}

  \item Roots are unique, i.e. for any node $u$ and $v$,
  \begin{align}
    \mathbf{root}[u] = \mathbf{root}[v] = 1\implies \mathbf{tree}[u]\neq \mathbf{tree}[v].
  \end{align}

  \item The consistency between parent and child, i.e.
  \begin{align}
    \mathbf{par}[\mathbf{lc}[u]]=\mathbf{par}[\mathbf{rc}[u]]=u,\\
    (\mathbf{lc}[\mathbf{par}[v]] - v)(\mathbf{rc}[\mathbf{par}[v]] - v) = 0,
  \end{align}
  for any valid internal node $u$ and valid non-root node $v$.
  Also, the parent of each root is dummy, i.e.
  \begin{align}
    \mathbf{par}\odot \mathbf{root} = \mathbf{0}.
  \end{align}

  \item The left and right children do not overlap, share the same tree index with the parent, and have correct depth,
  \begin{align}
    \mathbf{lc}[u] &\neq \mathbf{rc}[u],\\
  \mathbf{tree}[\mathbf{lc}[u]]&=\mathbf{tree}[\mathbf{rc}[u]]=\mathbf{tree}[u],\\
  \mathbf{depth}[\mathbf{lc}[u]]&=\mathbf{depth}[\mathbf{rc}[u]]=\mathbf{depth}[u]+1,
  \end{align}  
  for all internal nodes $u$.

  \item Leaves have no children, i.e. \begin{align}
    \mathbf{lc}\odot \mathbf{is\_leaf}=\mathbf{0}\land \mathbf{rc}\odot \mathbf{is\_leaf}=\mathbf{0},\label{eq:leaf-no-child}
  \end{align}
  where ``$\odot$'' is the Hadamard product.

  \item The hyper-box of any root is full, i.e.
  \begin{align}
    \mathbf{Root}\odot \mathbf{N}^{\leftarrow}=\mathbf{0}\land \mathbf{Root}\odot \mathbf{N}^{\rightarrow}=B\cdot \mathbf{Root},\label{eq:root-full-box}
  \end{align}
  where $\mathbf{Root} = (\mathbf{root}\|\cdots\|\mathbf{root})$.

  \item Valid nodes have non-empty hyper-box, i.e. for all node $u$ with non-zero depth,
    \begin{align}
    \mathbf{N}^{\rightarrow}[u, j] - \mathbf{N}^{\leftarrow}[u, j] > 0, \quad \forall j\in[d].
    \end{align}

  \item The hyper-box of left and right children are consistent with the parent and the split, i.e. for each internal node $u$ and feature $j$,
    \begin{align}
    \mathbf{N}^{\leftarrow}[u, j] - \mathbf{N}^{\leftarrow}[\mathbf{lc}[u], j] &= 0\land\label{eq:left-child-inherits-lowerbound}\\
    \mathbf{N}^{\rightarrow}[u, j] - \mathbf{N}^{\rightarrow}[\mathbf{rc}[u], j] &= 0\land\label{eq:right-child-inherits-upperbound}\\
    \mathbf{E}[u, j](\mathbf{N}^{\rightarrow}[\mathbf{lc}[u], j] - \mathbf{N}^{\leftarrow}[\mathbf{rc}[u], j]) &= 0\land\label{eq:split-child-consistency}\\
    (1-\mathbf{E}[u, j])(\mathbf{N}^{\leftarrow}[u, j] - \mathbf{N}^{\leftarrow}[\mathbf{rc}[u], j]) &= 0\land\label{eq:non-split-child-inherits-1}\\
    (1-\mathbf{E}[u, j])(\mathbf{N}^{\rightarrow}[u, j] - \mathbf{N}^{\rightarrow}[\mathbf{lc}[u], j]) &= 0.\label{eq:non-split-child-inherits-2}
    \end{align}
    \Cref{eq:left-child-inherits-lowerbound,eq:right-child-inherits-upperbound} ensure that 
    left child inherits the lower bound of the parent and right child inherits the upper bound.
    \Cref{eq:split-child-consistency} ensures that for the split feature, 
    the left child's upper bound meets the right child's lower bound.
    \Cref{eq:non-split-child-inherits-1,eq:non-split-child-inherits-2} together ensure that 
    for non-split features, the child hyper-box inherits the parent hyper-box.

  \item $\mathbf{E}$ is boolean and one-hot on internal nodes and zero on leaves, i.e. for each node $u$,
    \begin{align}
    \sum_{j\in[d]} \mathbf{E}[u, j] - \mathbf{valid}[u](1-\mathbf{is\_leaf}[u]) &= 0,\\
    \mathbf{E}\odot(\mathbf{E}-\mathbf{1}) &= \mathbf{0}.
    \end{align}

  \item $\mathbf{score}$ is zero on internal nodes, i.e.
    \begin{align}
    \mathbf{score}\odot (\mathbf{1}-\mathbf{is\_leaf}) = \mathbf{0}.
    \end{align}
  \end{enumerate}

The above involves essentially three kinds of techniques:
  proving routing, proving sum,
  and proving histogram.
  To prove $\mathbf{tree}[\mathbf{lc}[u]]=\mathbf{tree}[u]$ for all internal nodes $u$, the prover first
  commits $\mathbf{tree\_lc}$ where
  \[\mathbf{tree\_lc}[u] = \mathbf{tree}[\mathbf{lc}[u]]\]
  for all internal nodes $u$.
  It then proves
  \[
    \mathbf{lc}\|\mathbf{tree\_lc} \subseteq \left(\begin{matrix}
    0\\ 1\\ \vdots\\ N_{\mathrm{tot}}-1
    \end{matrix}\right)\|\mathbf{tree},
  \]
  alongside
  \[
    \mathbf{valid}\odot(\mathbf{1}-\mathbf{is\_leaf})\odot(\mathbf{tree\_lc} - \mathbf{tree}) = \mathbf{0}.
  \]

For sums over vectors, we use Aurora's sum-check protocol~\cite{aurora19}.
  To prove sum along one dimension of the matrix,
  we adopt a Quarks-style proof~\cite{quarks20},
  where the prover commits the prefix sum and total sum, 
  and proves the consistency between them and the original matrix.

Additionally, to prove the uniqueness of roots, we use
an approach that heralds our proof of histogram.
The prover commits a vector $\mathbf{cnt}\in \{0, 1\}^{TK+1}$, where $\mathbf{cnt}[t]$ counts the number of roots of tree $t$.
After proving its booleanity, the prover proves that
for a verifier challenge $\alpha$, 
\begin{align}
  \sum_{t=1}^{TK}\alpha^t = \sum_{t\in [TK+1]} \alpha^t\cdot \mathbf{cnt}[t] = \sum_{u\in [N_{\mathrm{tot}}]} \alpha^{\mathbf{tree}[u]}\cdot \mathbf{root}[u].
\end{align}
The LHS is public and can be computed efficiently.
To prove the $\alpha$-exponent, the prover commits a vector
that commits power-of-$\alpha$, and then proves the lookup relation.

\paragraph{Complexity.}
All forest-structural checks are low-degree algebraic constraints or
non-native predicates over either $N_{\mathrm{tot}}$ node rows or
$N_{\mathrm{tot}}d$ node-feature rows. The
only additional global object is the root-count vector of length $TK+1$.
Therefore the prover time for forest well-formedness is
\[
    \tilde{O}(N_{\mathrm{tot}}d + TK).
\]
Domain-lifting batching combines the zero-check and sum-check claims, and
interleaving batching combines the local lookup and range-check predicates.
As there are only constantly many polynomial commitments,
the verifier time and proof size are
$\tilde{O}(\log(N_{\mathrm{tot}}d+TK))$, up to the constant number of final KZG
openings.

\subsection{Proof of Training}\label{subsec:method-training}

The training proof reuses the same committed forest. 
The prover first commits the dataset $\mathbf{D}\in [B]^{M\times d}$ and 
its one-hot label $\mathbf{Y}\in \{0, 1\}^{K\times M}$.
The prover then commits leaves reached by each sample in the dataset,
by committing
\[
\mathbf{Ind} \in
\{N_{\mathrm{int}}+1,\ldots,N_{\mathrm{tot}}-1\}^{T\times K\times M},
\]
where $\mathbf{Ind}[t, k, r]$ 
is the index of the leaf reached by sample
$r$ in the tree trained for class $k$ at round $t$.

Likewise, the prover also commits the pseudo-residuals
$\mathbf{G}\in \mathbb{F}^{T\times K\times M}$ and the current scores
$\mathbf{S}\in \mathbb{F}^{T\times K\times M}$.
For round $t$, class $k$, and sample $r$,
$\mathbf{G}[t,k,r]$ represents the pseudo-residual target of sample $r$ that
should be approximated by the regression tree for class $k$ at round $t$, and
$\mathbf{S}[t,k,r]$ represents the cumulative score of sample $r$ for class
$k$ after round $t$ has been added.
Initially, $\mathbf{G}[0,k,r]=\mathbf{Y}[k,r]$ for all $k$ and $r$.
The committed leaf score stores the learning-rate-scaled tree contribution,
so the public learning rate $\mu$ is folded into the leaf-value relation.

Logically, the proof of training sequentially proves that:
\begin{enumerate}
  \item $\mathbf{Ind}$ correctly commits the leaf reached by each sample in the dataset.
    This is proved in a manner similar to proving batched inference.
  \item $\mathbf{G}[0, \cdot, \cdot]$ is set correctly from $\mathbf{Y}$,
    namely $\mathbf{G}[0,k,r]=\mathbf{Y}[k,r]$.
  \item The tree trained at each round $t$ is the optimal regression tree that fits $\mathbf{G}[t, \cdot, \cdot]$ on the dataset,
    i.e. each node split is optimal and the score is assigned to the leaves correctly. This is discussed in more detail in the next subsections.
  \item For each round $t$, the score $\mathbf{S}[t, \cdot, \cdot]$ is updated correctly
    from the previous score state and the scores each sample obtained from the trees trained at round $t$.
  \item For each round $t<T-1$, $\mathbf{G}[t+1, \cdot, \cdot]$ is updated correctly from $\mathbf{Y}$ and $\mathbf{S}[t, \cdot, \cdot]$.
    This involves proving softmax over the second dimension of $\mathbf{S}$.
\end{enumerate}

Together, these conditions imply that the committed objects encode a valid
plaintext GBDT training trace across all rounds.

\paragraph{Complexity.}
The training proof is dominated by three families of relations. First, routing
all training samples to their committed leaves costs
$\tilde{O}(TKMd)$. Second, the histogram proof processes the sample-update
stream of size $TKMd$ and the committed histogram tensors of size
$N_{\mathrm{tot}}dB$. Third, split optimality and leaf-value checks scan the
per-node, per-feature, and per-bin statistics, contributing
$\tilde{O}(N_{\mathrm{tot}}dB)$, while score and residual updates contribute
$\tilde{O}(TKM)$ plus the softmax-related fixed-point checks over the class
dimension. Overall, the prover time is
\[
    \tilde{O}\bigl(TKM(d+K) + N_{\mathrm{tot}}dB\bigr).
\]
The prover time is linear in the algebraic training trace, but the verifier
does not reprocess the trace: after batching, the verifier time and proof size
are
$\tilde{O}(\log(TKMd+N_{\mathrm{tot}}dB))$.

\subsection{Proving the Correctness of Node Split}\label{subsec:method-node-split}

Consider one node $u$, one feature $j$, and one threshold bin $b$. Let
$I_u \subseteq [M]$ be the samples reaching node $u$, and let
$L_{u,j,b},R_{u,j,b}$ be the induced left and right partitions. If each sample
$r\in I_u$ carries target value $g_r$, then the residual-sum-of-squares
objective is
\[
\operatorname{RSS}(u,j,b)
=
\sum_{r\in L_{u,j,b}}(g_r-\bar{g}_L)^2
+
\sum_{r\in R_{u,j,b}}(g_r-\bar{g}_R)^2.
\]
Expanding the squares gives
\begin{equation}\label{eq:rss-rewrite}
\operatorname{RSS}(u,j,b)
=
\sum_{r\in I_u} g_r^2
-
\left(
\frac{\left(\sum_{r\in L_{u,j,b}} g_r\right)^2}{|L_{u,j,b}|}
+
\frac{\left(\sum_{r\in R_{u,j,b}} g_r\right)^2}{|R_{u,j,b}|}
\right).
\end{equation}
For fixed node $u$, the first term in \eqref{eq:rss-rewrite} is independent of
the candidate split. Therefore minimizing $\operatorname{RSS}(u,j,b)$ is
equivalent to maximizing
\begin{equation}\label{eq:split-objective}
\Phi(u,j,b)
\coloneqq
\frac{\left(\sum_{r\in L_{u,j,b}} g_r\right)^2}{|L_{u,j,b}|}
+
\frac{\left(\sum_{r\in R_{u,j,b}} g_r\right)^2}{|R_{u,j,b}|},
\end{equation}
subject to the validity condition that both children are nonempty.

If we denote $\mathbf{C}[i, j, k]$ the number of samples falling to node $i$ with $j$-th feature being $k$,
  and $\mathbf{V}[i, j, k]$ the sum of residuals for samples falling to node $i$ with $j$-th feature being $k$,
  then we can rewrite \eqref{eq:split-objective} as
\begin{equation}
\Phi(u,j,b)
\coloneqq
\left\lfloor\frac{\left(\sum_{k=0}^{b-1} \mathbf{V}[u, j, k]\right)^2\cdot 2^q}{\sum_{k=0}^{b-1} \mathbf{C}[u, j, k]}\right\rfloor
+
\left\lfloor\frac{\left(\sum_{k=b}^{B-1} \mathbf{V}[u, j, k]\right)^2\cdot 2^q}{\sum_{k=b}^{B-1} \mathbf{C}[u, j, k]}\right\rfloor.
\end{equation}

This rewriting indicates that \emph{histograms} are sufficient proof objects. 
Once $\mathbf{C}$ and $\mathbf{V}$ are available, their prefix and suffix sums
let the prover prove the split gain $\Phi(u,j,b)$ for all candidate splits in one shot.\footnote{
  To avoid division by zero, we use an indicator
  that set the corresponding term to zero when the denominator is zero.
}

